%% ---------------------------------------------------------------------------------------------------------
%% Protocolo de evaluación del modelo de pronóstico — Big Data aplicado al fútbol
%%
%% Usa la clase pas.cls (la misma del protocolo de baloncesto), que ya trae a dos columnas, Times, natbib,
%% amsmath, booktabs, cleveref y hyperref: por eso aquí NO se carga ningún paquete extra.
%%
%% Si la clase se queja por `logo_cimat`, es la misma caja vacía que sale en el protocolo de baloncesto:
%% falta el archivo de imagen del logo, no es culpa de este documento.
%% ---------------------------------------------------------------------------------------------------------
\documentclass{pas}

\renewcommand\abstractname{Resumen}
\renewcommand\keywordname{Palabras clave}

\jnlDoiYr{2026}
\jnlPage{1}{10}
\doival{Borrador del 4 de octubre de 2026}

\articletitt{Protocolo de evaluación}
\title{Cómo juzgar un modelo de pronóstico de fútbol: protocolo de evaluación, referentes y un test
confirmatorio pre-registrado}

\author[César Carrillo]{\textrm}
\affil{Centro de Investigación en Matemáticas, unidad Aguascalientes}

\begin{document}

\begin{abstract}
Un modelo probabilístico de resultados de fútbol es fácil de evaluar mal: casi cualquier medida aislada admite una
lectura favorable. Este protocolo fija, antes de volver a mirar los números, qué se mide y contra qué. Separa
calidad de la probabilidad de acierto del pronóstico; calibración de capacidad de discriminación; y compara contra
un piso (frecuencias históricas) y un techo informativo (el mercado de apuestas con el margen removido), siempre
pareado partido a partido y con intervalos por bloques de jornada, porque los partidos de una misma jornada no son
independientes. Con 18{,}042 partidos fuera de muestra de siete ligas se establece que el modelo supera al piso y
pierde contra el mercado en las seis ligas con cuotas, por 0.008 a 0.010 de RPS, con intervalos que no cruzan el
cero. El hallazgo del proyecto —más empates en partidos parejos— se sometió a un test pre-registrado en Ligue~1 y
Eredivisie, las dos únicas ligas con cuotas que no se habían usado para encontrarlo: queda confirmado,
$+5.1$ puntos porcentuales $[+1.4, +8.7]$. La misma prueba muestra que ahí el modelo pierde contra el mercado
\emph{más} que en el resto, de modo que el patrón es real pero no es una ventaja. La distinción entre detectar una
regularidad y tener ventaja sobre quien ya la conoce es la razón de ser de este protocolo.
\end{abstract}

\keyword{pronóstico probabilístico, reglas de puntuación propias, calibración, pre-registro, mercados de apuestas,
bootstrap por bloques}

\maketitle

%% ===========================================================================================================
\section{Planteamiento}

Evaluar un pronóstico probabilístico involucra dimensiones que no se reducen una a otra:

\begin{itemize}
\item \textbf{Calidad de la probabilidad, no del acierto.} Un modelo que acierta más puede estar peor calibrado.
El acierto del \emph{pick} descarta lo que el modelo sí entrega: cuánta confianza tiene.
\item \textbf{Calibración frente a discriminación.} Son independientes: se puede estar perfectamente calibrado y no
distinguir un partido de otro, o discriminar bien y estar sistemáticamente sobreconfiado.
\item \textbf{Orden de las categorías.} Local, Empate y Visitante están ordenadas. Dar Visitante cuando fue Local
es peor que dar Empate.
\item \textbf{Dependencia entre observaciones.} Los partidos de una jornada comparten calendario y rachas; los de
un equipo, plantilla. Tratarlos como independientes subestima el error estándar.
\item \textbf{Existencia de un referente informado.} A diferencia de casi cualquier otro problema de predicción,
aquí hay un competidor público que agrega información que nosotros no tenemos: el mercado de apuestas.
\item \textbf{Multiplicidad y selección.} Siete ligas, varios subgrupos y varias métricas generan decenas de
comparaciones; sin control aparecerán hallazgos por azar.
\item \textbf{Hallazgos encontrados mirando los mismos datos.} Un hallazgo exploratorio no puede confirmarse con
los datos que lo produjeron, por bien que se mida.
\end{itemize}

%% ===========================================================================================================
\section{Etapas del análisis}

\subsection{Congelamiento del pronóstico}
Toda evaluación parte de pronósticos emitidos antes del partido y guardados con sello de tiempo. El ajuste
histórico usa validación \emph{walk-forward} por temporada: cada temporada de prueba se predice con un modelo
estimado solo con temporadas anteriores, y las variables pre-partido se construyen solo con partidos previos.

\textbf{Justificación.} Separa la reconstrucción del pasado del pronóstico. Solo lo segundo es evaluable.

\subsection{Calibración}
Se evalúa si las probabilidades significan lo que dicen mediante curvas de fiabilidad, calibración en el agregado y
error de calibración esperado, con una recalibración ajustada fuera de muestra para cuantificar cuánta mejora queda
sobre la mesa (figura~\ref{fig:calibracion}).

\textbf{Justificación.} Es la propiedad que vuelve utilizable un pronóstico para decidir.

\subsection{Puntuación propia y descomposición}
RPS como métrica principal, pérdida logarítmica y Brier, con la descomposición del Brier en fiabilidad, resolución
e incertidumbre.

\textbf{Justificación.} Las reglas propias no se pueden mejorar mintiendo sobre la propia confianza. La
descomposición dice si se pierde por calibración —corregible— o por no distinguir partidos.

\subsection{Comparación contra referentes}
Tres referentes en orden de exigencia: frecuencias históricas, mercado con el margen removido y, opcionalmente, un
Elo simple. La comparación es pareada y se acompaña de su intervalo por bloques.

\textbf{Justificación.} Una métrica sola no dice nada. Y la comparación debe ser pareada: la correlación entre
modelo y mercado reduce la desviación del diferencial de $0.55$ a $0.205$, que es lo que la hace viable.

\subsection{Valor de decisión}
Valor esperado a precios reales, criterios de apuesta y contraste contra el registro personal de jugadas.

\textbf{Justificación.} Un modelo puede ser peor en métricas y tener valor en un subconjunto, o al revés.

%% ===========================================================================================================
\section{Reglas de puntuación}

\subsection{Ranked Probability Score}
\begin{equation}
\mathrm{RPS} = \frac{1}{r-1}\sum_{i=1}^{r-1}\Bigg(\sum_{j\le i} p_j - \sum_{j\le i} o_j\Bigg)^{2},
\qquad r = 3
\end{equation}
con las categorías en orden Local $\prec$ Empate $\prec$ Visitante.

\textbf{Parámetros.} $p_j$, probabilidad pronosticada; $o_j$, indicador del resultado; $r$, número de categorías.

\textbf{Supuestos.} Que el orden es significativo, es decir, que el empate está entre las dos victorias.

\textbf{Utilidad.} Métrica estándar en fútbol: penaliza más los errores que cruzan el orden. Regla propia.

\textbf{Limitaciones.} Menos sensible que la logarítmica a probabilidades cercanas a cero; no descompone de forma
natural en calibración y resolución.

\subsection{Pérdida logarítmica}
\begin{equation}
\mathrm{LL} = -\frac{1}{n}\sum_{i=1}^{n}\log p_{i,y_i}
\end{equation}

\textbf{Utilidad.} Regla propia, interpretable en nats y ligada al crecimiento logarítmico del capital.

\textbf{Limitaciones.} No acotada: un solo partido con probabilidad casi nula domina el promedio.

\subsection{Brier y descomposición}
\begin{equation}
\mathrm{BS} = \frac{1}{n}\sum_{i=1}^{n}\sum_{j=1}^{r}(p_{ij}-o_{ij})^{2}
= \mathrm{fiab.} - \mathrm{resol.} + \mathrm{incert.}
\end{equation}

\textbf{Supuestos.} La descomposición es exacta solo con cubetas; su valor depende de cuántas, por lo que se
reporta la sensibilidad a esa elección.

%% ===========================================================================================================
\section{Referentes}

\subsection{Frecuencias históricas}
Probabilidad constante igual a las frecuencias de Local/Empate/Visitante del conjunto de entrenamiento. Es el
mínimo que cualquier modelo debe superar: no mirar el partido.

\subsection{Mercado}
\begin{equation}
p_i = \frac{1/c_i}{\sum_{j} 1/c_j}
\end{equation}

\textbf{Supuestos.} Que el margen de la casa se reparte proporcionalmente entre las tres categorías. Es una
simplificación conocida: la literatura documenta que el margen se carga más sobre los \emph{longshots}.

\textbf{Sensibilidad obligatoria.} Como el supuesto afecta al referente contra el que nos medimos, la comparación
se repite con los métodos de Shin y de \emph{odds-ratio}. Si la conclusión cambia según el método, la conclusión
es sobre el método y no sobre el modelo.

%% ===========================================================================================================
\section{Inferencia}

\subsection{Prueba pareada y bootstrap por bloques}
Para cada partido se calcula el diferencial de pérdida $d_i = S_i^{\mathrm{modelo}} - S_i^{\mathrm{ref}}$ y se
contrasta $\mathbb{E}[d]=0$. El intervalo se obtiene remuestreando \textbf{jornadas completas}, no partidos:
remuestrear partidos sueltos trataría como independiente lo que no lo es y daría intervalos demasiado estrechos.

\subsection{Multiplicidad}
Toda comparación secundaria se declara de antemano y se ajusta por Benjamini--Hochberg al 5\,\%. Las no declaradas
se reportan como exploratorias y no se interpretan como evidencia.

\subsection{Potencia}
El diferencial pareado de log-loss contra el mercado tiene desviación $\approx 0.205$. Para 80\,\% de potencia y
$\alpha = 0.05$ hacen falta los partidos del cuadro~\ref{tab:potencia}. Las siete ligas producen unos 32 partidos
por semana en temporada: reunir 450 lleva unas 14 semanas; reunir 3{,}300, unos dos años. Este cuadro es lo que
decide qué preguntas se pueden responder y cuáles hay que declarar fuera de alcance.

\begin{table}[t]
\caption{Partidos necesarios para detectar una diferencia dada contra el mercado, con 80\,\% de potencia.}
\label{tab:potencia}
\centering
\begin{tabular}{lr}
\hline
Diferencia (nats) & Partidos \\
\hline
0.030 & 367 \\
0.020 & 825 \\
0.010 & 3{,}299 \\
0.005 & 13{,}194 \\
\hline
\end{tabular}
\end{table}

%% ===========================================================================================================
\section{El test confirmatorio}
\label{sec:test}

El proyecto observó que en partidos con probabilidades parejas —rango máximo menos mínimo $\le 0.125$— el empate es
más frecuente. Sobre las cinco ligas donde se exploró, los parejos son el 15.2\,\% de 13{,}283 partidos y en ellos
el empate ocurre el 28.5\,\% de las veces frente al 24.5\,\% en el resto: $+3.9$ puntos $[+1.9, +6.1]$. El hallazgo
se encontró explorando esa misma muestra, así que se trató como hipótesis a confirmar.

\subsection{Pre-registro}
Las reglas se fijaron y se subieron al repositorio \textbf{antes de descargar los datos}; el historial de control de
versiones es la prueba. Se fijó: la población (Ligue~1 y Eredivisie completas, las dos únicas ligas con cuotas no
usadas para encontrar el hallazgo); el umbral de $0.125$, sin probar ningún otro; la especificación del modelo, con
ventana expandible por ser la de las otras ligas de 18 equipos y \textbf{sin probar alternativas}; la medida
primaria; la inferencia por bloques de jornada; y la regla de decisión en sus tres casos, incluida la de retirar la
afirmación si el intervalo incluía el cero.

\subsection{Resultado}
Sobre 4{,}759 partidos fuera de muestra, 785 parejos (16.5\,\%), el empate ocurre el 29.2\,\% en los parejos frente
al 24.1\,\% en el resto: \textbf{$+5.1$ puntos $[+1.4, +8.7]$}. El intervalo excluye el cero con signo positivo,
luego \textbf{la hipótesis queda confirmada en muestra independiente}. El efecto es incluso algo mayor que en la
muestra que la generó, lo que descarta que fuera un artefacto de selección (figura~\ref{fig:parejos}).

\begin{figure*}[t]
\centering
\includegraphics[width=\textwidth]{figuras/parejos}
\caption{Diferencia en la frecuencia de empate entre partidos parejos y el resto. En azul la muestra
confirmatoria, pre-registrada. Barras: intervalo al 95\,\% por bootstrap de bloques de jornada.}
\label{fig:parejos}
\end{figure*}

\subsection{La lectura que corrige al proyecto}
La segunda medida secundaria declarada —RPS del modelo contra el mercado restringido a los parejos— da
$+0.0093$ $[+0.0065, +0.0120]$ frente a $+0.0083$ $[+0.0072, +0.0094]$ en el resto. En los partidos parejos el
modelo pierde contra el mercado \emph{algo más} que en el resto, no menos. Y las dos fuentes asignan prácticamente
la misma probabilidad al empate ahí —modelo 28.3\,\%, mercado 28.0\,\%, tasa real 29.2\,\%—: el mercado conoce el
patrón igual de bien.

Así que «los partidos parejos empatan más» es cierto y ahora está confirmado; «ahí es donde el modelo aporta algo»
es falso. Es donde el modelo describe bien una regularidad que el mercado describe igual de bien, y donde sigue
perdiendo (figura~\ref{fig:mercado}).

\begin{figure*}[t]
\centering
\includegraphics[width=\textwidth]{figuras/mercado}
\caption{Diferencia de RPS entre el modelo y el mercado, por liga y por subconjunto. A la derecha del cero gana el
mercado. Ningún intervalo cruza el cero.}
\label{fig:mercado}
\end{figure*}

%% ===========================================================================================================
\section{Estado actual del modelo}

Con 18{,}042 partidos fuera de muestra de siete ligas, de los cuales 16{,}615 tienen cuotas, el modelo supera
claramente al piso de frecuencias y pierde contra el mercado en las seis ligas con cuotas
(cuadro~\ref{tab:estado}). La desventaja no es ruido y la muestra está holgadamente sobrada para detectarla: se
necesitaban entre 370 y 650 partidos y hay entre 1{,}900 y 3{,}100 por liga.

\begin{table}[t]
\caption{Pérdida logarítmica fuera de muestra. El piso son las frecuencias históricas.}
\label{tab:estado}
\centering
\begin{tabular}{lrrrr}
\hline
Liga & $n$ & Modelo & Mercado & Piso \\
\hline
Liga MX        & 1{,}427 & 1.030 & ---   & 1.065 \\
Premier League & 3{,}114 & 0.981 & 0.957 & 1.068 \\
La Liga        & 3{,}141 & 0.999 & 0.973 & 1.069 \\
Bundesliga     & 2{,}543 & 1.007 & 0.979 & 1.075 \\
Serie A        & 3{,}058 & 0.981 & 0.956 & 1.082 \\
Ligue 1        & 2{,}837 & 1.011 & 0.983 & 1.076 \\
Eredivisie     & 1{,}922 & 0.966 & 0.937 & 1.072 \\
\hline
\end{tabular}
\end{table}

\begin{figure}[t]
\centering
\includegraphics[width=\columnwidth]{figuras/calibracion}
\caption{Calibración de la probabilidad de empate. Los dos pronósticos van ligeramente por encima de la diagonal
en la parte baja: declaran menos empates de los que ocurren.}
\label{fig:calibracion}
\end{figure}

\subsection{Limitación de fondo}
El esquema \emph{walk-forward} evita la fuga temporal en los coeficientes, pero la elección de las variables, de la
ventana y de la familia del modelo se hizo mirando el conjunto completo. Eso es selección sobre toda la muestra, un
nivel por encima del que controla la validación. La única corrección real es evaluar sobre datos posteriores a esa
elección, que es lo que acumula el seguimiento congelado.

%% ===========================================================================================================
\section{Conclusión metodológica}

El uso combinado de estas piezas permite separar la calidad de la probabilidad de la del acierto; distinguir fallo
de calibración de fallo de información; medir contra un piso y contra un techo informativo en vez de en abstracto;
cuantificar qué preguntas admite el tamaño de muestra; y evitar que un hallazgo exploratorio se publique como
confirmado. El caso de los partidos parejos muestra las dos caras: el protocolo confirmó el patrón y, con la misma
prueba, desmintió la interpretación con la que se venía publicando.

\end{document}
